Euler and Pressure

For inviscid flow:
$$\nabla p=\rho(\vec g-\vec a)$$
Horizontal x-y plane: gx=gy=0:
$$p_x=-\rho a_x$$
$$p_y=-\rho a_y$$
Vertical y upward: gy = -g:
$$p_y=\rho(-g-a_y)$$

EXAMPLE FROM ACCELERATION NOTE
rho = 1000 kg/m^3; a=(-7,14).
Horizontal plane: px=+7000 Pa/m; py=-14000 Pa/m.
If y vertical UP: py=-23810 Pa/m.
Always draw gravity before selecting signs.

PRESSURE FIELD FROM GRADIENT
Integrate px with respect to x and add F(y,z,t). Determine F using py and pz. A reference p is required for the final additive constant.

Example steady field:
$u=Ax+By$, $v=Bx-Ay$.
Let $C=A^2+B^2$; a=(Cx,Cy).
Horizontal plane:
$$p=p_{ref}-\frac{\rho C}{2}S$$
$$S=x^2-x_{ref}^2+y^2-y_{ref}^2$$

WITH VISCOSITY
Do not use Euler when the viscous term matters:
$$\nabla p=\rho\vec g+\mu\nabla^2\vec V-\rho\vec a$$
